\section{Experiments}
\label{sec:exp}

We evaluate how much ViBridge-VQA improves Vintern-1B, how many tokens it should use,
where its gains come from, and whether they carry over to unseen datasets.

\subsection{Dataset and Evaluation Metrics}

\paragraph{Dataset.}
We evaluate on AutoViVQA~\cite{autovivqa2026}, a Vietnamese VQA benchmark built on MS COCO
images~\cite{lin2015coco} with an LLM-driven generation pipeline and ensemble-based
quality control. Each question has five short free-form reference answers of 1--10 words
and is labelled with its reasoning type. Because the public random 80/20 split lets questions about the same image appear in
both training and validation, we remove 273 duplicated (image, question) pairs and 97 questions
with irregular type labels, keep the eight remaining reasoning categories, and re-partition
the 36{,}707 questions by image into a 70/15/15 split (seed 42). The resulting training,
validation and test sets contain 25{,}776, 5{,}463 and 5{,}468 questions on 13{,}579,
2{,}909 and 2{,}910 images, respectively. This grouped split is slightly harder than a
random one (by at most 0.5 F1 in our runs), and all our results are reported on it.

\paragraph{Evaluation metrics.}
Following AutoViVQA and ViMoE-VQA, we report answer-matching metrics (exact-match
Accuracy, and word-level Precision, Recall and F1) and generation metrics
(BLEU~\cite{papineni2002bleu}, ROUGE-L~\cite{lin2004rouge},
METEOR~\cite{banerjee2005meteor} and CIDEr~\cite{vedantam2015cider}), all multiplied by
100. Word-level scores are computed against each of the five references, keeping the best
match. Note that we average F1 over questions, which yields lower values than the harmonic
mean of the averaged precision and recall used for the published baselines; computed in
the latter way, ViBridge-VQA would obtain 55.95 instead of 54.56 F1. Our F1 comparisons
are therefore conservative.

\subsection{Experimental Setup}

\paragraph{Implementation details.}
The vision encoder, the projector and the decoder are initialised from
Vintern-1B-v3.5~\cite{doan2024vintern1b} and kept frozen. ViBridge-VQA is trained for two
epochs, after which neither the validation loss nor CIDEr improves further; the variants
with a decoder LoRA are trained for one or three epochs. We use AdamW with a learning rate
of $2\times10^{-4}$ and an effective batch size of 8, train on the first reference answer,
and decode greedily. Unless stated otherwise, every configuration is trained with three random
seeds (42, 123 and 3407), and we report the mean and the standard deviation. All runs fit on a single 16\,GB
GPU (Tesla P100 or T4). The token budget is selected on the validation split, and the test
split is used only for the final report.

\paragraph{Baselines.}
We compare against three groups of systems: (1) multimodal models fine-tuned on
AutoViVQA, namely ViT5\_ViT,
BARTPhoBEiT~\cite{tran2023bartphobeitpretrainedsequencetosequenceimage}, Vintern-1B and
ViMoE-VQA~\cite{vimoevqa}; (2) models used zero-shot, namely Vintern-1B, LLaMA~3.2,
Gemini~2.0 and 2.5 Flash, and GPT-5; and (3) our \emph{global-only} model, which uses the
same bridge without local tokens ($k=0$, $g=8$), with and without a decoder LoRA. The
results of groups (1) and (2) are taken from the AutoViVQA and ViMoE-VQA papers and were
obtained on an 8:1:1 split of the original release. The fine-tuned Vintern-1B follows the
official recipe: the vision encoder and the projector are frozen, a rank-16 LoRA adapts the
language model for one epoch, and each image is encoded as up to six $448\times448$ tiles
plus a thumbnail.

\subsection{Main Results}
\label{sec:main}

\begin{table}[t]
\centering
\caption{Performance comparison on the AutoViVQA test set ($\times100$). Baselines are
taken from AutoViVQA~\cite{autovivqa2026} and ViMoE-VQA~\cite{vimoevqa}; ViMoE-VQA is
averaged over five seeds. For our models we report the mean $\pm$ standard deviation over
three seeds (global only: four), the share of trained parameters, and use 8 (global only),
50 ($k{=}36$) or 158 ($k{=}144$) visual tokens, against up to 1{,}792 for the fine-tuned
Vintern-1B. Best results per column are in bold. $^\dagger$Far outside the range of all
other systems and not used for comparison.}
\label{tab:main}
\footnotesize
\setlength{\tabcolsep}{3pt}
\resizebox{\textwidth}{!}{%
\begin{tabular}{llrrrrrrrr}
\toprule
Category & Model & Acc & Prec & Rec & F1 & BLEU & ROUGE-L & METEOR & CIDEr \\
\midrule
\multirow{4}{*}{\shortstack[l]{Multimodal\\(fine-tuned)}}
 & ViT5\_ViT & 7.97 & 46.84 & 50.33 & 48.52 & 4.13 & 46.89 & 31.02 & 72.68 \\
 & BARTPhoBEiT & 8.81 & 45.30 & 46.48 & 45.88 & 43.29$^\dagger$ & 44.83 & 24.57 & 188.96$^\dagger$ \\
 & Vintern-1B & 13.01 & 52.47 & 55.12 & 53.76 & 6.11 & 51.93 & 35.25 & 72.84 \\
 & ViMoE-VQA & 9.55 & \textbf{62.89} & 58.60 & \textbf{60.69} & 12.54 & 47.12 & 39.10 & 88.76 \\
\midrule
\multirow{5}{*}{\shortstack[l]{LLMs and VLMs\\(zero-shot)}}
 & Vintern-1B & 0.12 & 17.52 & 19.87 & 17.55 & 1.91 & 25.84 & 23.93 & 8.54 \\
 & LLaMA 3.2 & 0.36 & 23.96 & 73.71 & 36.16 & 3.62 & 36.11 & 30.01 & 62.84 \\
 & Gemini 2.0 Flash & 0.55 & 27.20 & 74.10 & 39.79 & 4.41 & 39.60 & 31.72 & 74.42 \\
 & Gemini 2.5 Flash & 0.22 & 24.43 & \textbf{76.66} & 24.75 & 0.39 & 37.27 & 31.22 & 71.90 \\
 & GPT-5 & 10.84 & 47.20 & 55.20 & 50.89 & 6.07 & 47.30 & 33.34 & 84.20 \\
\midrule
\multirow{2}{*}{\shortstack[l]{Global only\\(ours, $k{=}0$)}}
 & $g{=}8$ (0.78\%) & \pms{8.86}{.37} & \pms{51.19}{.19} & \pms{52.37}{.21} & \pms{50.47}{.22} & \pms{16.24}{.43} & \pms{48.70}{.19} & \pms{41.10}{.29} & \pms{96.19}{1.10} \\
 & + LoRA, 3 ep.\ (1.01\%) & \pms{11.40}{.24} & \pms{55.37}{.18} & \pms{56.16}{.09} & \pms{54.52}{.14} & \pms{20.97}{.29} & \pms{52.68}{.21} & \pms{45.08}{.14} & \pms{108.43}{.67} \\
\midrule
\multirow{3}{*}{\shortstack[l]{ViBridge-VQA\\(ours, $g{=}14$)}}
 & $k{=}36$ (1.35\%) & \pms{11.75}{.21} & \pms{55.36}{.21} & \pms{56.56}{.21} & \pms{54.56}{.19} & \pms{18.52}{.18} & \pms{52.64}{.19} & \pms{45.53}{.33} & \pms{110.58}{.95} \\
 & $k{=}144$ (1.35\%) & \pms{13.19}{.23} & \pms{56.64}{.13} & \pms{58.17}{.26} & \pms{56.04}{.19} & \pms{19.79}{.15} & \pms{54.11}{.16} & \pms{47.38}{.26} & \pms{116.74}{.76} \\
 & $k{=}36$ + LoRA, 1 ep.\ (1.57\%) & \pms{\textbf{14.67}}{.30} & \pms{59.02}{.12} & \pms{59.96}{.33} & \pms{58.17}{.19} & \pms{\textbf{22.53}}{.11} & \pms{\textbf{56.25}}{.22} & \pms{\textbf{49.32}}{.37} & \pms{\textbf{122.88}}{.93} \\
\bottomrule
\end{tabular}}
\end{table}

Table~\ref{tab:main} compares ViBridge-VQA with the baselines. Without updating any
weight of Vintern-1B, ViBridge-VQA raises its F1 from 17.55 in the zero-shot setting to
54.56, and its CIDEr from 8.54 to 110.58. With the selected budget of 50 visual tokens, it
already outperforms the officially fine-tuned Vintern-1B on seven of the eight metrics; the
largest gains appear in the generation metrics (BLEU, METEOR and CIDEr), and only
exact-match accuracy remains 1.3 points lower. When all 144 local tokens are used, the
model surpasses the fine-tuned Vintern-1B on every metric, including accuracy, while still
reading up to eleven times fewer visual tokens. These results suggest that much of
what AutoViVQA requires is already present in the frozen Vintern-1B, provided that its
decoder receives suitable visual input.

Decoder adaptation adds to these gains rather than replacing them. Training ViBridge-VQA
jointly with a rank-16 LoRA for a single epoch raises test F1 by a further 3.6 points to
58.17 and gives the best accuracy, BLEU, ROUGE-L, METEOR and CIDEr among all compared systems
(setting aside the outlying BARTPhoBEiT values). The task-specific ViMoE-VQA, whose fusion module and decoder are trained on
AutoViVQA, keeps the highest precision and F1, 2.5 F1 points above our best model. Zero-shot language models reach a high recall with long answers but lose precision,
whereas our answers match the reference length (Section~\ref{sec:analysis}).

Our global-only model isolates the effect of the local tokens: even with a frozen decoder, ViBridge-VQA matches the
global-only model trained with a three-epoch decoder LoRA in F1 (54.56 versus 54.52) and
exceeds it in CIDEr (110.58 versus 108.43), although its BLEU is lower. 

\paragraph{Statistical stability.}
All our models behave consistently across random seeds. On the test split, the standard
deviation of ViBridge-VQA is at most 0.19 for F1 and 0.95 for CIDEr, and for each of our models in
Table~\ref{tab:main}, test F1 lies at most 0.6 points below validation F1, which indicates that the
budget selected on the validation split does not over-fit it.

\subsection{Choosing the Token Budget}
\label{sec:budget-results}

\begin{table}[t]
\centering
\caption{Effect of the token budget ($\times100$, mean over three seeds, $\pm$ standard
deviation for F1; global only: four seeds). Tokens: visual tokens passed to the decoder
($g+k$). Time: minutes for one evaluation session on the 10{,}931 validation and test
questions, including model loading (Kaggle P100/T4). The selected budget is in bold.}
\label{tab:budget}
\scriptsize
\setlength{\tabcolsep}{4.5pt}
\begin{tabular}{rrrrrrrr}
\toprule
 & & & & \multicolumn{2}{c}{Validation} & \multicolumn{2}{c}{Test} \\
\cmidrule(lr){5-6}\cmidrule(lr){7-8}
$g$ & $k$ & Tokens & Time & F1 & CIDEr & F1 & CIDEr \\
\midrule
8 & 0 & 8 & -- & \pms{50.74}{.17} & 99.07 & \pms{50.47}{.22} & 96.19 \\
\midrule
8 & 1 & 9 & 57.0 & \pms{51.03}{.12} & 99.27 & \pms{50.71}{.16} & 96.90 \\
8 & 9 & 17 & 59.2 & \pms{52.82}{.11} & 106.11 & \pms{52.11}{.25} & 101.30 \\
8 & 36 & 44 & 60.6 & \pms{54.32}{.41} & 111.88 & \pms{53.82}{.22} & 107.67 \\
8 & 144 & 152 & 71.2 & \pms{56.15}{.08} & 119.05 & \pms{55.20}{.07} & 114.13 \\
\midrule
14 & 1 & 15 & 58.8 & \pms{51.67}{.25} & 101.66 & \pms{51.67}{.12} & 99.53 \\
14 & 9 & 23 & 57.7 & \pms{52.98}{.08} & 106.02 & \pms{52.76}{.28} & 103.63 \\
\textbf{14} & \textbf{36} & \textbf{50} & \textbf{60.0} & \pms{\textbf{54.96}}{.28} & \textbf{113.29} & \pms{\textbf{54.56}}{.19} & \textbf{110.58} \\
14 & 144 & 158 & 70.2 & \pms{56.61}{.04} & 120.71 & \pms{56.04}{.19} & 116.74 \\
\bottomrule
\end{tabular}
\end{table}

\begin{figure}[t]
\centering
\includegraphics[width=\linewidth]{fig_results}
\caption{(a) Validation F1 against the number of visual tokens (mean $\pm$ std over three
seeds); the circled point is the knee, $g=14$, $k=36$. Horizontal lines: the global-only
model (four seeds) and the same model with a three-epoch decoder LoRA (three seeds).
(b) Validation F1 by reasoning type (number of questions in parentheses): global only and
+ LoRA for seed 42, ViBridge-VQA ($g=14$, $k=36$) without and with a one-epoch decoder LoRA
averaged over three seeds; the numbers give the gain of ViBridge-VQA + LoRA over the
global-only model.}
\label{fig:results}
\end{figure}

Table~\ref{tab:budget} and Fig.~\ref{fig:results}a vary the number of global tokens
$g\in\{8,14\}$ and of local tokens $k\in\{1,9,36,144\}$. The most striking pattern is that
\emph{local tokens matter far more than global ones}. At $g=14$, increasing $k$ from 1 to
144 adds 4.9 F1 and 19.1 CIDEr on the validation split, whereas increasing $g$ from 8 to 14
at a fixed $k$ adds only 0.2--0.6 F1. Accuracy has not saturated within the 144 region
tokens available at our resolution, as each step of $k$ still adds 1.3--2.0 F1, but the gain
per token drops quickly, from 0.07 F1 per token between 23 and 50 tokens to 0.015 between
50 and 158. The budget is therefore a genuine trade-off between accuracy and cost.

Because accuracy grows with every added token, all configurations lie on the Pareto front,
and we select its knee (Section~\ref{sec:budget}). On the validation split, the knee is
$g=14$, $k=36$. Remarkably, the same configuration is the knee for each of the eight
metrics on both validation and test, whether cost is counted in visual tokens or measured
as inference time; with a logarithmic token axis, it remains the knee in 11 of the 16
cases. TOPSIS leads to the same conclusion when accuracy and cost are weighted
equally, using the gains in F1, CIDEr and METEOR as benefits and the added tokens,
inference time, trained parameters and seed variance as costs. Weighting cost more favours
$k=9$, and weighting accuracy more favours $k=144$, but $g=14$, $k=36$ is the only
configuration that stays among the top four for every accuracy weight between 0.3 and 0.7.
In practical terms, it obtains 72\% of the F1 gain of $k=144$ over the global-only
model with 32\% of its visual tokens, and it is the cheapest configuration that reaches the global-only model with a
three-epoch decoder LoRA. Fourteen global tokens are better than eight at every $k$, so we
use $g=14$.

\paragraph{Computational efficiency.}
ViBridge-VQA encodes a single $336\times336$ view per image, which costs 362 GFLOPs and
229\,ms with InternViT-300M on a Tesla P100. The six $448\times448$ tiles and the thumbnail
used by the fine-tuned Vintern-1B cost 2{,}172 GFLOPs and 922\,ms, and give the decoder up to
1{,}792 visual tokens instead of 50. Keeping all 144 local tokens instead of 36 makes evaluation about 17\% slower.

\subsection{Ablation Study}
\label{sec:ablation}



\begin{table}[t]
\centering
\caption{Where the visual tokens come from, with and without decoder adaptation
(validation, $\times100$, mean over three seeds with $\pm$ std for F1; Multi-Token: four
seeds; $^\ast$seed 42 only). Bridge-only models are trained for two epochs; + LoRA is a
rank-16 decoder LoRA trained jointly with the bridge for one epoch (2.16M additional
parameters). Params: trained bridge parameters.}
\label{tab:bridges}
\scriptsize
\setlength{\tabcolsep}{4pt}
\begin{tabular}{lrrrrrr}
\toprule
 & & & \multicolumn{2}{c}{Bridge only} & \multicolumn{2}{c}{+ decoder LoRA} \\
\cmidrule(lr){4-5}\cmidrule(lr){6-7}
Visual tokens from & Tokens & Params & F1 & CIDEr & F1 & CIDEr \\
\midrule
\multicolumn{7}{l}{\emph{Global vector (\texttt{[CLS]})}} \\
\quad Residual & 1 & 4.86M & \pms{46.02}{.36} & 86.48 & \pms{52.69}{.14} & 104.25 \\
\quad Multi-Token & 8 & 7.35M & \pms{50.74}{.17} & 99.07 & \pms{53.47}{.17} & 106.68 \\
\multicolumn{7}{l}{\emph{Learned compression of the patch tokens}} \\
\quad Tile-Attention & 8 & 4.14M & \pms{45.14}{.92} & 84.14 & 52.82$^\ast$ & 104.72$^\ast$ \\
\quad Light Q-Former & 8 & 27.6M & \pms{47.06}{.66} & 88.07 & \pms{53.21}{.14} & 106.26 \\
\quad Full Q-Former & 16 & 69.4M & \pms{46.84}{.47} & 87.74 & \pms{53.29}{.29} & 106.50 \\
\midrule
\multicolumn{7}{l}{\emph{Global vector + pre-trained projector (ViBridge-VQA, $g=14$)}} \\
\quad $k=36$ & 50 & 12.86M & \pms{54.96}{.28} & 113.29 & \pms{\textbf{58.21}}{.16} & \textbf{124.88} \\
\quad $k=144$ & 158 & 12.86M & \pms{\textbf{56.61}}{.04} & \textbf{120.71} & -- & -- \\
\bottomrule
\end{tabular}
\end{table}

\paragraph{Where should the visual tokens come from?}
Table~\ref{tab:bridges} compares three sources of visual tokens for the frozen decoder. The
first group reads only the global \texttt{[CLS]} vector, through a \emph{Residual} bridge
that produces one token or the \emph{Multi-Token} bridge of Eq.~\eqref{eq:bridge}. The
second group learns to compress the patch tokens, either with a \emph{Tile-Attention} bridge
in which eight learned queries attend over the patches, or with a light or full
Q-Former~\cite{li2023blip2}. Without decoder adaptation, the Multi-Token bridge is the best
of these five, and the bridges that compress the patches are 3.7--5.6 F1 lower despite
having up to $9.4\times$ more parameters. Learning a new alignment of the patch tokens from
25k questions thus does not pay off. ViBridge-VQA, which takes the local tokens from the
pre-trained projector instead, improves over the best global-only bridge by 4.2 F1 with
$k=36$ and by 5.9 F1 with $k=144$. The same pooling operation therefore helps when it is
applied to the output of the pre-trained projector, whereas in a preliminary study, pooling
the raw patch tokens through a newly learned projection lost 5.1--13.9 F1 and a
Honeybee-style convolutional abstractor~\cite{cha2024honeybee} remained 0.7 F1 below the
Multi-Token bridge with $2.7\times$ its parameters.

\paragraph{Local tokens and decoder adaptation are complementary.}
A decoder LoRA improves every global-only bridge and brings them within 0.8 F1 of one
another (52.69--53.47), which suggests that the decoder learns to compensate for poor visual
input. Even after this adaptation, all of them remain below ViBridge-VQA with a frozen
decoder. Decoder adaptation is therefore not needed to surpass them, but it is still
useful: in the same one-epoch joint training, it adds 3.3 F1 to ViBridge-VQA, more than the
2.7 F1 it adds to the Multi-Token bridge. Conversely, distilling the projector into
the global-only bridge, at the feature or output level, lowers F1 by 0.3--0.4 points: the
projector's information is best passed as tokens.

\subsection{Analysis by Reasoning Type and Error Analysis}
\label{sec:analysis}

\paragraph{Reasoning types.}
Fig.~\ref{fig:results}b breaks the results down by the reasoning categories of AutoViVQA.
Local tokens help most where the answer depends on what is visible in the image:
recognition (+10.2 F1 over the global-only model), action (+5.4) and spatial (+5.1)
questions. On these categories, ViBridge-VQA with a frozen decoder is also better than the
global-only model with a decoder LoRA. Conversely, decoder adaptation is more useful for
causal and relational questions, whose reference answers are the longest (5.0 and 5.6 words
on average), and for yes/no questions. Adding the LoRA to ViBridge-VQA combines both
strengths: the resulting model is the best in seven of the eight categories, matching the
global-only model with LoRA on causal and relational questions while keeping the large gain
of the local tokens on recognition (58.7 versus 49.7 F1). Local tokens improve what the model
perceives, whereas decoder adaptation improves how it formulates reasoning answers.

\paragraph{Error analysis.}
An analysis of our validation predictions helps explain the low exact-match accuracy of
all systems. Our answers are as long as the references on average (4.3
words), so the gains do not come from changing the answer length. Most predictions (77--81\%) share some, but not all, words with the closest reference: a
correct answer is often phrased differently from all five references. Local tokens turn some of these
partial answers into exact ones and reduce complete failures: compared with the global-only model,
ViBridge-VQA increases the share of predictions with F1 = 1 from 9.8\% to 13.3\% and reduces
the share with no word in common with any reference from 9.6\% to 8.0\% (6.8\% with LoRA).
The remaining failures concentrate on context, recognition, spatial and action questions
(12--18\% without any overlap), i.e., on fine-grained perception, whereas they are rare for
counting, causal and relational questions ($<5\%$), where a partially correct answer is
easier to produce.

\subsection{Generalisation to Other Datasets}
\label{sec:ood}

\begin{table}[t]
\centering
\caption{Zero-shot transfer to Vietnamese VQA datasets not seen during training
($\times100$; about 1{,}000 sampled test questions per sampling seed, 505--519 for ViVQA;
mean over three sampling seeds, $\pm$ std for F1). Vintern-1B reads six tiles and a
thumbnail; the other models are trained on AutoViVQA only (seed 42). Best results per
column are in bold.}
\label{tab:ood}
\scriptsize
\setlength{\tabcolsep}{2.5pt}
\resizebox{\textwidth}{!}{%
\begin{tabular}{lrrrrrrrr}
\toprule
 & \multicolumn{2}{c}{ViTextVQA} & \multicolumn{2}{c}{OpenViVQA} & \multicolumn{2}{c}{ViVQA-X} & \multicolumn{2}{c}{ViVQA} \\
\cmidrule(lr){2-3}\cmidrule(lr){4-5}\cmidrule(lr){6-7}\cmidrule(lr){8-9}
Model & F1 & CIDEr & F1 & CIDEr & F1 & CIDEr & F1 & CIDEr \\
\midrule
Vintern-1B, zero-shot & \pms{32.01}{.31} & \textbf{195.62} & \pms{\textbf{58.40}}{.80} & \textbf{411.63} & \pms{14.52}{.47} & 50.13 & \pms{23.90}{.24} & 78.63 \\
Global only + LoRA, 3 ep. & \pms{3.96}{.09} & 15.64 & \pms{21.56}{.48} & 83.21 & \pms{15.89}{.65} & 42.12 & \pms{31.48}{1.15} & 84.94 \\
\midrule
ViBridge-VQA, $k{=}36$ & \pms{15.14}{.47} & 66.39 & \pms{28.24}{.79} & 128.36 & \pms{24.91}{.14} & 79.51 & \pms{42.68}{1.52} & 171.21 \\
ViBridge-VQA, $k{=}144$ & \pms{\textbf{33.26}}{.67} & 190.30 & \pms{35.89}{.80} & 187.00 & \pms{\textbf{27.89}}{.13} & 92.22 & \pms{\textbf{44.95}}{.90} & \textbf{191.28} \\
ViBridge-VQA, $k{=}36$ + LoRA & \pms{15.39}{.39} & 70.07 & \pms{28.99}{.51} & 130.46 & \pms{27.64}{.11} & \textbf{95.29} & \pms{44.76}{1.21} & 183.60 \\
\bottomrule
\end{tabular}}
\end{table}

Table~\ref{tab:ood} evaluates the models without further training on four Vietnamese VQA
datasets: two scene-text datasets, ViTextVQA~\cite{nguyen2025vitextvqa} and
OpenViVQA~\cite{nguyen2023openvivqa}, and two general-domain ones,
ViVQA-X~\cite{vivqax2025} and ViVQA~\cite{tran2021vivqa}. We include the strongest
global-only model, trained with a three-epoch decoder LoRA, as a reference. Local tokens
change the picture most on scene text, where the global-only model collapses, presumably
because the \texttt{[CLS]} vector cannot convey printed words. On ViTextVQA, F1 rises from 3.96 to 15.14
with 36 local tokens and to 33.26 with 144, slightly above the zero-shot Vintern-1B that
reads six $448\times448$ tiles, although ViBridge-VQA reads a single $336\times336$ view. On
OpenViVQA, whose references are long descriptive sentences, the local tokens narrow the gap
to Vintern-1B, but Vintern-1B remains clearly better.

On the general-domain datasets, ViBridge-VQA outperforms both reference models, by 10.4
(ViVQA-X) and 18.8 (ViVQA) F1 over the zero-shot Vintern-1B with $k=36$. As on AutoViVQA,
the zero-shot model reaches a high recall with long answers (61.2 on ViVQA-X) but a very low
precision (9.1). Adding the decoder LoRA to ViBridge-VQA helps on these general-domain
datasets, bringing $k=36$ close to $k=144$, but barely on scene text. Reading text thus
requires more local tokens, which decoder adaptation cannot replace, consistent with the
token-budget results above.

\subsection{Qualitative Analysis}

\begin{figure}[!tb]
\centering
\includegraphics[width=\linewidth]{fig_qualitative}
\caption{Qualitative examples from the AutoViVQA validation set (seed 42). Q: question,
GT: one of the five reference answers, Global: global-only model, +LoRA: the same model with
a three-epoch decoder LoRA, Ours: ViBridge-VQA ($g=14$, $k=36$). Plausible answers are shown
in green and wrong ones in red.}
\label{fig:qual}
\end{figure}

In Examples 1--3 of Fig.~\ref{fig:qual}, only ViBridge-VQA answers correctly, and each
answer depends on a local detail: the brand printed on the
laptop, the canal that identifies Venice, and the water sprayed by the elephant. The
global-only models instead answer ``Apple'', ``Hanoi'' and ``swimming'', which are plausible
for the scene as a whole but wrong for the image. Example~4 shows a typical failure:
ViBridge-VQA describes the scene instead of explaining the purpose of the bus, which both
global-only models infer. This matches the advantage of decoder adaptation on causal questions.
